\section{Implications for lattice calculations}
\label{sec:lattice}

Our proof in \sec{ope} makes clear the relationship between the renormalized quasi-PDF, \spcorr, and pseduo-PDF distributions.  As a practical matter there are a few different ways in which these equations can be used to convert a lattice calculation of the \spcorr $\tilde Q$ into a PDF. Three examples are 1) first Fourier transform to the quasi-PDF with \eq{fac-tq-orig0}, and then use \eq{fac-tq}, 2) first Fourier transform to the pseudo-PDF with \eq{pseudoRen}, and then use \eq{ps-q-fact}, and 3) first match to the Fourier transform the position space PDF $Q(\zeta,\mu)$ using \eq{io-Q-fact}, and then transform it to the PDF with the inverse of \eq{FTq}.  Since the numerical implementation of these steps may have slightly different systematics it is interesting to compare them, or to use more than one approach in order to reduce uncertainties.

According to the analysis in Sec.~\ref{sec:ope}, for the factorization formula of the Euclidean distributions to work, one must calculate the same \spcorr with small distance $z^2$ and large momentum $P^z$ so that the dynamical and kinematic higher-twist effects are suppressed. For practical lattice calculations, this means that there is only a finite number of useful data points in $(z, P^z)$ that we can use to extract the PDF.

To illustrate this, consider a $48^3\times64$ lattice with spacing $a=0.09$ fm. The distance of the spatial correlation $z$ is in units of $a\sim1/2.2\ {\rm GeV}^{-1}$, and the nucleon momentum $P^z$ is in units of $2\pi/(48a)\sim0.29$ GeV. Let us take $\Lambda_{\rm QCD}\sim 0.3$ GeV. In principal the target mass corrections can be subtracted. If we consider various values $z=ma$ and $P^z=n*2\pi/(48a)$ for integer $m$ and $n$, then to satisfy $z\Lambda_{\rm QCD}\ll1$ and $P^z\gg \Lambda_{\rm QCD}$, we must have
\beq
 m\ll 11\,,\,\,\,\quad n\gg 1\,.
\eeq
To control the higher-twist correction at $20\%$, i.e. $z^2\Lambda_{\rm QCD}^2\sim0.2,\ \Lambda_{\rm QCD}^2/P_z^2\sim0.2$, we can only choose
\beq
m = \{0,1,2,3,4\}\,,\,\,\,\, n=\{3,4,5,6,\cdots\}\,,
\eeq
where the largest value for $n$ is limited by what is practical in current lattice simulations. Six is the largest number of units attained in Ref.~\cite{Orginos:2017kos}. For quasi-PDF calculations, there are $4\times2+1=9$ useful data points for each fixed momentum $|P^z|$; for pseudo-PDF calculation, there are only $4\times2=8$ useful data points for each fixed $|z|$. In either case, it is anticipated that a direct Fourier transform with respect to $z$ or $\zeta=zP^z$ will lead to oscillation in $x$-space and incorrect prediction for the small-$x$ region due to the truncation in coordinate space. This has been observed in a recent lattice calculation of the quasi-PDF in Ref.~\cite{Chen:2017mzz}. Methods have been developed in recent works to eliminate the oscillation from the truncation effect~\cite{Lin:2017ani,Zhang:2017gau} in the quasi-PDF, while the higher-twist contributions at large $z$ still need to be systematically corrected. It should be noted that the above is a rough estimate of the higher-twist corrections since the prefactor of $z^2\Lambda_{\rm QCD}^2$ could be smaller than 1. Their actual significance can only be quantitatively determined from lattice simulations.

To fully take advantage of all the useful data points, we can evolve them to either the same $z^2$ or $P^z$ according to the perturbative analysis, which has been studied in Refs.~\cite{Orginos:2017kos,Karpie:2017bzm,Radyushkin:2017lvu} for the \spcorr. However, since the evolution equation of the \spcorr in $\ln z^2$ or $\ln P_z^2$ follows a nonlocal convolution in $\zeta=zP^z$ or $z$, one has to know the full information in coordinate space to do the evolution. With limited number of data points, either large uncertainties or adopting a model-dependent assumption about the shape is inevitable.

To improve the precision of either approach, the only way forward is to have finer lattice spacing $a$ so that we could have more data points which satisfy $|z|\ll \Lambda_{\rm QCD}^{-1}$ and larger nucleon momentum $P^z$.
With increasing $P^z$, the valence distribution of the nucleon is contracted in the $z$ direction, so the spatial correlation of valence quarks is shrinked into smaller distance in $z$. If $P^z$ is large enough, the spatial correlation will fall off quickly within $|z|< \Lambda_{\rm QCD}^{-1}$, then the truncation error from Fourier transform will be significantly reduced. On the other hand, if we interpret the spatial correlation as the \spcorr, its shape will not change under a Lorentz boost because it is a scalar function of $\zeta=zP^z$ and $z^2$. Nevertheless, finer lattice spacing $a$ allows for calculation with a wider range of $P^z$, thus covering larger values of $\zeta=zP^z$ to reduce the truncation error. Since the number of useful data points increases quadratically with $1/a$, a more precise lattice calculation with controlled systematic errors will be available in the future.
